% !TEX root = ../master.tex

\section{RecursiveSet: API und fachliche Modellierung}
\label{sec:recursive-set-api}

Nachdem im vorigen Kapitel die typseitigen Grundlagen von TypeScript
eingeführt wurden, wendet sich die Arbeit nun dem zentralen technischen
Baustein der Portierung zu: \mintinline{text}{RecursiveSet}. Dabei geht es
zunächst noch nicht um die interne Implementierung, sondern um die Frage,
welches fachliche Verhalten diese Datenstruktur bereitstellt und warum sie für
die Übertragung der Lehrnotebooks notwendig wurde.

Der Ausgangspunkt ist dabei ein sehr konkretes Problem. In den ursprünglichen
Python-Notebooks wird mit Mengen so gearbeitet, wie man es aus der
mathematischen Sicht erwarten würde. Python beschreibt \mintinline{python}{set}
als ungeordnete Sammlung ohne Duplikate und nennt ausdrücklich
Mitgliedschaftstests sowie Operationen wie Vereinigung, Schnitt, Differenz und
symmetrische Differenz als typische Einsatzfälle \autocite{python_datastructures}.
Genau dieses Verhalten wird in den Lehrmaterialien an vielen Stellen
vorausgesetzt.

Bei einer direkten Übertragung nach TypeScript tritt jedoch ein Bruch auf.
Zwar besitzt auch JavaScript mit \mintinline{text}{Set} eine eingebaute
Mengenstruktur, diese vergleicht Objekte jedoch nicht nach ihrem Inhalt,
sondern nach ihrer Identität, also danach, ob es sich um exakt dieselbe
Referenz handelt \autocite{mdn_set}. Für primitive Werte fällt dieser
Unterschied kaum auf. Bei Tupeln, Mengen von Tupeln oder Mengen von Mengen
wird er jedoch zum zentralen Problem der gesamten Portierung.

Dieses Kapitel erklärt deshalb \mintinline{text}{RecursiveSet} bewusst aus
Anwendersicht. Im Mittelpunkt steht nicht, wie Hashing, Einfügen oder Löschen
intern funktionieren, sondern wie sich die Datenstruktur im Code benutzen
lässt und wie sie das aus Python bekannte Mengenverhalten wiederherstellt.
Die technische Umsetzung folgt erst im nächsten Kapitel.

\subsection{Ziel des Kapitels}
\label{subsec:rs-api-ziel}

Ziel dieses Kapitels ist es, \mintinline{text}{RecursiveSet} als fachliche
Arbeitsoberfläche für Mengen in der TypeScript-Portierung einzuführen. Die
Datenstruktur soll nicht einfach eine weitere Sammlung neben Array,
\mintinline{text}{Set} oder Map sein. Ihre eigentliche Aufgabe besteht darin,
die aus den Python-Notebooks vertraute Denkweise in Mengen für TypeScript
wieder verfügbar zu machen.

Für die Leserinnen und Leser ist dabei ein Perspektivwechsel wichtig. In
Kapitel~2 stand im Vordergrund, wie TypeScript Typen beschreibt und wie
dadurch Programmfehler früher sichtbar werden. In diesem Kapitel geht es nun
um eine konkrete Datenstruktur, die auf diesen Typideen aufbaut und ein sehr
spezifisches fachliches Problem löst: zusammengesetzte Werte sollen in Mengen
so behandelt werden, wie es in der theoretischen Informatik erwartet wird.

Das Kapitel verfolgt daher drei eng zusammenhängende Teilziele. Erstens wird
gezeigt, an welcher Stelle sich das eingebaute \mintinline{text}{Set} für
unsere Zwecke unerwartet verhält. Zweitens wird \mintinline{text}{RecursiveSet}
als gezielte Antwort auf dieses Problem eingeführt. Drittens wird damit die
Grundlage für die späteren API-Abschnitte gelegt, in denen Mengenoperationen,
Transformationen und verschachtelte Mengen schrittweise erklärt werden.

Didaktisch ist dabei besonders wichtig, zunächst beim beobachtbaren Verhalten
zu bleiben. Leserinnen und Leser sollen zuerst verstehen, \emph{welches}
Problem auftritt und \emph{welches} Verhalten \mintinline{text}{RecursiveSet}
bereitstellt. Erst danach ist es sinnvoll zu untersuchen, \emph{wie} dieses
Verhalten intern technisch erreicht wird.

\subsection{Das Mengenproblem beim Übergang von Python nach TypeScript}
\label{subsec:rs-api-mengenproblem}

In den Python-Notebooks wirken Mengen zunächst vollkommen unproblematisch.
Python beschreibt Mengen als Datenstruktur für eindeutige Elemente,
Mitgliedschaftstests und mathematische Mengenoperationen \autocite{python_datastructures}.
Gerade für die theoretische Informatik ist das äußerst passend, weil Zustände,
Übergänge oder Symbolmengen direkt als Mengen modelliert werden können.

Ein einfaches Beispiel mit Tupeln zeigt dieses erwartete Verhalten:

\begin{listing}[H]
	\caption{Mengen in Python: Das erwartete Verhalten bei Tupeln}
	\label{lst:set_python_tuple}
	\begin{minted}{python}
		menge = set()
		menge.add((1, 2))
		
		print(menge)           # {(1, 2)}
		print((1, 2) in menge) # True
	\end{minted}
\end{listing}

Das Ergebnis entspricht hier genau der mathematischen Intuition. Das Tupel
\mintinline{python}{(1, 2)} wurde eingefügt und wird bei einer späteren
Mitgliedschaftsabfrage auch wieder erkannt. Für die Arbeit mit Mengen in den
Lehrnotebooks ist dieses Verhalten so selbstverständlich, dass es meist gar
nicht eigens thematisiert werden muss.

Bei der Übertragung nach TypeScript entsteht an genau dieser Stelle jedoch ein
unerwarteter Bruch. Das eingebaute JavaScript-\mintinline{text}{Set} speichert
zwar ebenfalls nur eindeutige Werte, bei Objekten basiert Gleichheit jedoch
auf Identität, also auf derselben Referenz, nicht auf gleichem Inhalt
\autocite{mdn_set}. Für zusammengesetzte Werte führt das zu einem Ergebnis,
das aus fachlicher Sicht überraschend wirkt:

\begin{listing}[H]
	\caption{Mengen in TypeScript: Das fehlerhafte Verhalten bei Tupeln}
	\label{lst:set_ts}
	\begin{minted}{typescript}
		// Wir definieren den Typ Tuple, der zwei Zahlen akzeptiert.
		type Tuple = [number, number];
		
		// Die Menge soll nur diese Tupel akzeptieren
		const menge = new Set<Tuple>();
		menge.add([1, 2]);
		
		// Die Menge enthält das Element sichtbar
		console.log(menge);             // Ausgabe: Set(1) { [ 1, 2 ] }
		
		// Die Abfrage liefert ein ungewöhnliches Ergebnis
		console.log(menge.has([1, 2])); // Ausgabe: false
	\end{minted}
\end{listing}

Gerade dieses Beispiel ist für den weiteren Verlauf der Arbeit entscheidend.
Das Problem besteht nicht darin, dass \mintinline{text}{Set} \enquote{falsch}
implementiert wäre. Das Problem besteht vielmehr darin, dass sich das
Standardverhalten von JavaScript-\mintinline{text}{Set} nicht mit der
fachlichen Erwartung deckt, die in den Python-Notebooks stillschweigend
vorausgesetzt wird.

Für einfache Demo-Programme mag dieser Unterschied noch harmlos erscheinen.
In Algorithmen der theoretischen Informatik hat er jedoch unmittelbare Folgen.
Wenn Mengen von Zuständen, Zustandspaaren oder Teilmengen verwendet werden,
dann hängt die Korrektheit ganzer Berechnungen davon ab, dass logisch gleiche
Werte in der Menge auch wirklich als gleich erkannt werden. Geschieht das
nicht, schlagen Mitgliedschaftsabfragen fehl, Duplikate werden nicht erkannt
und algorithmische Invarianten gehen verloren.

Genau an dieser Stelle zeigt sich also, dass die Portierung nicht durch eine
reine Syntaxübersetzung gelöst werden kann. Die Sprache TypeScript ist nicht
das eigentliche Problem. Entscheidend ist vielmehr, dass ihre eingebaute
Mengenstruktur für zusammengesetzte Werte ein anderes Gleichheitsverständnis
besitzt als dasjenige, auf dem die Python-Notebooks fachlich aufbauen.

\subsection{RecursiveSet als Mengenkompensation}
\label{subsec:rs-api-kompensation}

Um diesen Bruch zu beheben, wurde \mintinline{text}{RecursiveSet} entwickelt.
Die Datenstruktur verfolgt nicht das Ziel, das eingebaute
\mintinline{text}{Set} vollständig zu ersetzen. Ihre Aufgabe ist viel
spezifischer: Sie soll für die Portierung genau das Mengenverhalten
bereitstellen, das in den ursprünglichen Python-Notebooks bereits vorausgesetzt
wird.

Man kann \mintinline{text}{RecursiveSet} daher am besten als
\enquote{Mengenkompensation} verstehen. Der Begriff meint hier nicht bloß eine
kleine technische Korrektur, sondern den bewussten Ausgleich eines
sprachbedingten Unterschieds. Python und JavaScript stellen beide eine
Mengenstruktur bereit, aber sie behandeln zusammengesetzte Werte nicht in
derselben Weise. \mintinline{text}{RecursiveSet} gleicht genau diesen
Unterschied aus.

Für die Nutzung im Projekt bedeutet das: Wenn zwei Werte fachlich dieselbe
Struktur besitzen, dann sollen sie auch in einer Menge als gleich gelten.
Gerade bei Tupeln, Mengen von Tupeln oder Mengen von Mengen ist das die
entscheidende Eigenschaft. Nicht die zufällige Identität eines konkreten
Objekts im Speicher soll maßgeblich sein, sondern der Inhalt, den dieser Wert
repräsentiert.

Das lässt sich bereits an einem einfachen Gegenbeispiel zum vorherigen
\mintinline{text}{Set}-Beispiel andeuten:

\begin{listing}[H]
	\caption{\texttt{RecursiveSet} als fachlich passende Mengenstruktur}
	\label{lst:recursive_set_tuple_intro}
	\begin{minted}{typescript}
		const menge = new RecursiveSet<Tuple<[number, number]>>();
		menge.add(new Tuple(1, 2));
		
		console.log(menge.has(new Tuple(1, 2))); // true
	\end{minted}
\end{listing}

An diesem Beispiel ist zunächst nur die fachliche Wirkung wichtig. Obwohl bei
der Abfrage ein neues Objekt erzeugt wird, erkennt die Menge, dass es sich
inhaltlich um denselben Wert handelt. Genau dieses Verhalten macht
\mintinline{text}{RecursiveSet} für die Portierung der Lehrnotebooks so
zentral.

Damit ist zugleich die Perspektive für die nächsten Abschnitte festgelegt.
Im weiteren Verlauf von Kapitel~3 wird nicht mehr gefragt, ob
\mintinline{text}{RecursiveSet} \enquote{auch eine Collection} ist. Stattdessen
wird untersucht, wie sich mit dieser Datenstruktur die aus Python vertrauten
Mengenoperationen, verschachtelten Mengen und projektbezogenen Algorithmen
wieder präzise formulieren lassen. Erst wenn dieses API-Verhalten klar ist,
lohnt sich der Blick auf die interne technische Umsetzung in Kapitel~4.

\subsection{Grundlegende Benutzung im Vergleich zu Python}
\label{subsec:rs-api-grundbenutzung}

Nachdem im vorigen Abschnitt geklärt wurde, warum das eingebaute
\mintinline{text}{Set} für unsere Portierung nicht ausreicht, kann nun die
eigentliche Benutzung von \mintinline{text}{RecursiveSet} betrachtet werden.
Dabei ist der Vergleich zu Python besonders wichtig, weil die Lehrnotebooks
bereits eine sehr natürliche und fachlich passende Arbeitsweise mit Mengen
etabliert haben.

Python beschreibt Mengen als Datenstruktur für eindeutige Elemente und hebt
insbesondere Mitgliedschaftstests sowie typische Mengenoperationen hervor
\autocite{python_datastructures}. Genau an dieser vertrauten Arbeitsweise
orientiert sich auch \mintinline{text}{RecursiveSet}. Ziel ist also nicht, eine
möglichst originelle neue API zu entwerfen, sondern eine Benutzung
bereitzustellen, die für die Portierung fachlich denselben Zweck erfüllt.

Ein erster Vergleich zeigt bereits, dass sich die grundlegenden Operationen
sehr ähnlich lesen lassen. In Python wird eine Menge erzeugt, anschließend
werden Elemente hinzugefügt, entfernt und auf Mitgliedschaft geprüft:

\begin{listing}[H]
	\caption{Grundlegende Mengenbenutzung in Python}
	\label{lst:python_set_basic_usage}
	\begin{minted}{python}
		numbers = set()
		
		numbers.add(1)
		numbers.add(2)
		numbers.add(2)
		
		print(numbers)        # {1, 2}
		print(1 in numbers)   # True
		
		numbers.remove(1)
		print(numbers)        # {2}
		print(len(numbers))   # 1
	\end{minted}
\end{listing}

Dasselbe Grundmuster lässt sich mit \mintinline{text}{RecursiveSet} in
TypeScript formulieren:

\begin{listing}[H]
	\caption{Grundlegende Mengenbenutzung mit \texttt{RecursiveSet}}
	\label{lst:ts_recursive_set_basic_usage}
	\begin{minted}{typescript}
		const numbers = new RecursiveSet<number>();
		
		numbers.add(1);
		numbers.add(2);
		numbers.add(2);
		
		console.log(numbers);        // {1, 2}
		console.log(numbers.has(1)); // true
		
		numbers.remove(1);
		console.log(numbers);        // {2}
		console.log(numbers.size);   // 1
	\end{minted}
\end{listing}

Gerade an diesem einfachen Beispiel ist die eigentliche Stärke von
\mintinline{text}{RecursiveSet} noch gar nicht spektakulär. Bei primitiven
Werten wie Zahlen verhält sich die Datenstruktur zunächst genau so, wie man es
von einer Menge erwartet. Das ist didaktisch sogar ein Vorteil: Die API soll
an dieser Stelle bewusst vertraut wirken und keine unnötige zusätzliche
Komplexität einführen.

Zugleich werden hier bereits die wichtigsten Grundoperationen sichtbar.
\mintinline{text}{new RecursiveSet(...)} übernimmt die Rolle der Konstruktion,
\mintinline{text}{add} fügt ein Element ein, \mintinline{text}{remove} entfernt
ein Element, \mintinline{text}{has} prüft die Mitgliedschaft und
\mintinline{text}{size} liefert die Anzahl der gespeicherten Werte. Damit
entsteht eine Arbeitsoberfläche, die aus Sicht der Notebook-Portierung sehr
nahe an der bisherigen Python-Nutzung liegt.

Für die Lehrnotebooks ist diese Nähe entscheidend. Studierende sollen beim
Lesen eines portierten Algorithmus nicht zuerst eine vollkommen neue
Collections-Bibliothek erlernen müssen. Vielmehr soll die Mengenidee aus der
theoretischen Informatik im Vordergrund bleiben, während sich die konkrete
Benutzung in TypeScript möglichst natürlich anfühlt.

Der eigentliche Unterschied zu Python zeigt sich daher noch nicht bei Zahlen
oder Strings, sondern erst bei zusammengesetzten Werten. Genau dort wird
sichtbar, warum \mintinline{text}{RecursiveSet} mehr ist als nur eine weitere
Mengenklasse mit bekannten Methoden. Dieser Punkt wird im folgenden Abschnitt
schrittweise aufgebaut.

\subsection{Gleichheit und Mitgliedschaft bei zusammengesetzten Werten}
\label{subsec:rs-api-gleichheit}

Die grundlegende Benutzung mit primitiven Werten ist wichtig, aber sie erklärt
noch nicht, warum \mintinline{text}{RecursiveSet} für die Portierung überhaupt
notwendig wurde. Der entscheidende Unterschied zeigt sich erst bei
zusammengesetzten Werten, also bei Werten, die selbst wieder aus mehreren
Bestandteilen bestehen. Gerade solche Strukturen treten in den Lehrnotebooks
ständig auf, etwa als Zustandspaare, Kanten, Übergänge oder Mengen von
Zuständen.

In Python wirkt diese Art der Mengenbenutzung sehr natürlich. Ein Tupel kann
direkt in einer Menge gespeichert und später über eine logisch gleiche
Tupelstruktur wiedergefunden werden:

\begin{listing}[H]
	\caption{Mitgliedschaft bei zusammengesetzten Werten in Python}
	\label{lst:python_set_tuple_membership}
	\begin{minted}{python}
		pairs = set()
		pairs.add((1, 2))
		
		print((1, 2) in pairs)  # True
	\end{minted}
\end{listing}

Gerade diese Kürze macht die Python-Lösung so elegant. Ein Tupel kann direkt
als zusammengesetzter Wert verwendet werden, und die Mengenabfrage orientiert
sich an der fachlichen Struktur des Werts. Für die Notebooks ist das äußerst
angenehm, weil die mathematische Idee fast ohne zusätzlichen technischen
Aufwand im Code erscheint.

Mit \mintinline{text}{RecursiveSet} lässt sich dieselbe Grundidee auch in
TypeScript wieder präzise ausdrücken. Dabei übernimmt die Klasse
\mintinline{text}{Tuple} die Rolle eines strukturell vergleichbaren
zusammengesetzten Werts:

\begin{listing}[H]
	\caption{Mitgliedschaft bei zusammengesetzten Werten mit \texttt{RecursiveSet}}
	\label{lst:ts_recursive_set_tuple_membership}
	\begin{minted}{typescript}
		const pairs = new RecursiveSet<Tuple<[number, number]>>();
		pairs.add(new Tuple(1, 2));
		
		console.log(pairs.has(new Tuple(1, 2))); // true
	\end{minted}
\end{listing}

Dieses Beispiel ist inhaltlich sehr wichtig. Bei der Abfrage wird nicht
dieselbe Instanz wiederverwendet, sondern ein neues Objekt erzeugt. Trotzdem
liefert die Mitgliedschaftsabfrage \mintinline{text}{true}. Genau daran zeigt
sich, dass für \mintinline{text}{RecursiveSet} nicht die zufällige Identität
eines Objekts entscheidend ist, sondern seine fachliche Struktur.

An dieser Stelle kann nun auch der zentrale Begriff dieses Kapitels eingeführt
werden: \textit{Wertsemantik}. Damit ist gemeint, dass zwei Werte dann als
gleich behandelt werden, wenn sie inhaltlich dieselbe Struktur besitzen. Für
unsere Portierung ist genau dieses Verhalten entscheidend. Ein Zustandspaar
soll nicht nur dann erkannt werden, wenn zufällig exakt dasselbe Objekt erneut
verwendet wird, sondern auch dann, wenn an anderer Stelle ein logisch gleiches
Paar neu erzeugt wurde.

Der Unterschied wird besonders deutlich, wenn man mehrere logisch gleiche
Werte in dieselbe Menge einfügt. In Python entsteht daraus weiterhin nur ein
einziger Mengeneintrag:

\begin{listing}[H]
	\caption{Python erkennt logisch gleiche Tupel als gleich}
	\label{lst:python_set_tuple_dedup}
	\begin{minted}{python}
		pairs = set()
		pairs.add((1, 2))
		pairs.add((1, 2))
		
		print(pairs)      # {(1, 2)}
		print(len(pairs)) # 1
	\end{minted}
\end{listing}

Dasselbe Verhalten stellt \mintinline{text}{RecursiveSet} auch in TypeScript
wieder her:

\begin{listing}[H]
	\caption{\texttt{RecursiveSet} erkennt logisch gleiche Tupel als gleich}
	\label{lst:ts_recursive_set_tuple_dedup}
	\begin{minted}{typescript}
		const pairs = new RecursiveSet<Tuple<[number, number]>>();
		pairs.add(new Tuple(1, 2));
		pairs.add(new Tuple(1, 2));
		
		console.log(pairs.toString()); // {(1, 2)}
		console.log(pairs.size);       // 1
	\end{minted}
\end{listing}

Damit wird aus einer bloß ähnlichen API nun tatsächlich eine fachlich passende
Mengenstruktur. \mintinline{text}{RecursiveSet} übernimmt also nicht nur
Methodennamen wie \mintinline{text}{add} oder \mintinline{text}{has}, sondern
stellt genau das Gleichheitsverständnis bereit, das für die Python-Notebooks
bereits selbstverständlich war.

Für die theoretische Informatik ist das von zentraler Bedeutung. Viele
Algorithmen arbeiten mit Mengen von Zuständen, Paaren oder Teilmengen. Wenn
logisch gleiche Werte dabei nicht zuverlässig als gleich erkannt würden,
wären Duplikaterkennung, Mitgliedschaftsabfragen und zahlreiche
Korrektheitsargumente gefährdet. Erst durch diese strukturbezogene
Gleichheit wird die Mengenmodellierung in TypeScript wieder fachlich belastbar.

Zugleich bereitet dieser Abschnitt den nächsten Schritt vor. Wenn schon ein
einzelnes Zustandspaar als zusammengesetzter Wert sinnvoll in einer Menge
gespeichert werden kann, dann liegt der Übergang zu eigentlichen
Mengenoperationen nahe. Im folgenden Abschnitt wird daher gezeigt, wie
\mintinline{text}{RecursiveSet} die aus Python vertrauten Operationen wie
Vereinigung, Schnitt und Differenz wieder als direkte Arbeitsmittel verfügbar
macht.


\subsection{Klassische Mengenoperationen im direkten Vergleich}
\label{subsec:rs-api-mengenoperationen}

Nachdem die grundlegende Benutzung von \mintinline{text}{RecursiveSet}
eingeführt wurde, kann nun der eigentliche fachliche Kern betrachtet werden:
die klassischen Mengenoperationen. Gerade an dieser Stelle war die Arbeit mit
Python in den ursprünglichen Lehrnotebooks besonders elegant. Die
Python-Dokumentation nennt für Mengen ausdrücklich die mathematischen
Operationen Vereinigung, Schnitt, Differenz und symmetrische Differenz
\autocite{python_datastructures}. Genau diese Denkweise sollte in der
TypeScript-Portierung erhalten bleiben.

Für dieses Kapitel ist dabei ein didaktischer Punkt wichtig. Es geht nicht
darum, die Syntax von Python möglichst exakt nachzuahmen. Entscheidend ist
vielmehr, dass dieselben fachlichen Aussagen weiterhin direkt im Code
ausgedrückt werden können. \mintinline{text}{RecursiveSet} soll also nicht wie
Python aussehen, sondern sich bei Mengenoperationen fachlich so verhalten,
dass dieselben mathematischen Überlegungen weiterhin natürlich formulierbar
bleiben.

\paragraph*{Vereinigung}
Die Vereinigung zweier Mengen enthält alle Elemente, die in mindestens einer
der beiden Mengen vorkommen. In Python lässt sich diese Idee besonders knapp
mit dem Operator \mintinline{python}{|} ausdrücken; die Python-Dokumentation
zeigt genau diese Schreibweise als Standardbeispiel für Mengenoperationen
\autocite{python_datastructures}.

\begin{listing}[H]
	\caption{Vereinigung in Python und mit \texttt{RecursiveSet}}
	\label{lst:union_python_ts}
	\begin{minted}{python}
		A = {1, 2, 3}
		B = {3, 4, 5}
		
		print(A | B)   # {1, 2, 3, 4, 5}
	\end{minted}
	
	\vspace{0.5em}
	
	\begin{minted}{typescript}
		const A = new RecursiveSet(1, 2, 3);
		const B = new RecursiveSet(3, 4, 5);
		
		console.log(A.union(B)); // {1, 2, 3, 4, 5}
	\end{minted}
\end{listing}

Der Unterschied liegt hier vor allem in der Oberflächensyntax. Python benutzt
einen sehr kompakten Operator, während \mintinline{text}{RecursiveSet} eine
explizite Methode bereitstellt. Fachlich bleibt die Aussage jedoch dieselbe:
Das Ergebnis enthält alle Werte, die in mindestens einer der beiden Mengen
enthalten sind.

\paragraph*{Schnitt}
Der Schnitt zweier Mengen enthält genau die Werte, die in beiden Mengen
zugleich vorkommen. Auch hierfür verwendet Python eine knappe Operatorsyntax,
nämlich \mintinline{python}{&} \autocite{python_datastructures}.

\begin{listing}[H]
	\caption{Schnitt in Python und mit \texttt{RecursiveSet}}
	\label{lst:intersection_python_ts}
	\begin{minted}{python}
		A = {1, 2, 3}
		B = {2, 3, 4}
		
		print(A & B)   # {2, 3}
	\end{minted}
	
	\vspace{0.5em}
	
	\begin{minted}{typescript}
		const A = new RecursiveSet(1, 2, 3);
		const B = new RecursiveSet(2, 3, 4);
		
		console.log(A.intersection(B)); // {2, 3}
	\end{minted}
\end{listing}

Gerade diese Operation ist für die theoretische Informatik besonders wichtig,
weil sie häufig beim Vergleich von Zustandsmengen, Symbolmengen oder
Ergebnismengen verwendet wird. Für die Portierung war daher entscheidend, dass
der Schnitt nicht nur für primitive Werte, sondern auch für zusammengesetzte
Werte fachlich zuverlässig funktioniert.

\paragraph*{Differenz}
Die Differenz \(A \setminus B\) enthält alle Elemente, die in \(A\), aber nicht
in \(B\) vorkommen. Die Python-Dokumentation demonstriert diese Operation mit
dem Operator \mintinline{python}{-} \autocite{python_datastructures}.

\begin{listing}[H]
	\caption{Differenz in Python und mit \texttt{RecursiveSet}}
	\label{lst:difference_python_ts}
	\begin{minted}{python}
		A = {1, 2, 3}
		B = {2, 4}
		
		print(A - B)   # {1, 3}
	\end{minted}
	
	\vspace{0.5em}
	
	\begin{minted}{typescript}
		const A = new RecursiveSet(1, 2, 3);
		const B = new RecursiveSet(2, 4);
		
		console.log(A.difference(B)); // {1, 3}
	\end{minted}
\end{listing}

Auch hier zeigt sich sehr schön, was mit Mengenkompensation gemeint ist. Die
Schreibweise ist in TypeScript etwas ausführlicher, aber die fachliche Idee
bleibt vollständig erhalten. Wer die Mengendifferenz aus Python oder aus der
Mathematik kennt, erkennt dieselbe Operation unmittelbar wieder.

\paragraph*{Symmetrische Differenz}
Die symmetrische Differenz enthält alle Werte, die in genau einer der beiden
Mengen vorkommen, also nicht in beiden gleichzeitig. Python verwendet dafür
den Operator \mintinline{python}{^}; auch dieses Beispiel erscheint direkt in
der Python-Dokumentation zu Mengen \autocite{python_datastructures}.

\begin{listing}[H]
	\caption{Symmetrische Differenz in Python und mit \texttt{RecursiveSet}}
	\label{lst:symmetric_difference_python_ts}
	\begin{minted}{python}
		A = {1, 2, 3}
		B = {3, 4, 5}
		
		print(A ^ B)   # {1, 2, 4, 5}
	\end{minted}
	
	\vspace{0.5em}
	
	\begin{minted}{typescript}
		const A = new RecursiveSet(1, 2, 3);
		const B = new RecursiveSet(3, 4, 5);
		
		console.log(A.symmetricDifference(B)); // {1, 2, 4, 5}
	\end{minted}
\end{listing}

Gerade diese Operation wirkt im ersten Moment oft etwas weniger vertraut als
Vereinigung oder Schnitt. Für viele Algorithmen ist sie jedoch äußerst
nützlich, weil sie genau die Werte sichtbar macht, bei denen sich zwei Mengen
unterscheiden.

\paragraph*{Warum dieser Vergleich für die Portierung so wichtig ist}
An den vier Beispielen wird sehr deutlich, weshalb der Vergleich mit Python in
diesem Kapitel im Vordergrund steht. Die ursprünglichen Lehrnotebooks konnten
Mengenoperationen äußerst knapp und fachlich direkt formulieren, weil Python
für Mengen genau diese mathematischen Grundoperationen bereits als natürliche
Sprachmittel bereitstellt \autocite{python_datastructures}.

\mintinline{text}{RecursiveSet} übernimmt diese Eleganz nicht auf Ebene der
Syntax, wohl aber auf Ebene der fachlichen Modellierung. Die Operationen
heißen nun \mintinline{text}{union}, \mintinline{text}{intersection},
\mintinline{text}{difference} und \mintinline{text}{symmetricDifference}, aber
sie stellen genau dieselben mathematischen Zusammenhänge bereit. Für die
Portierung der Lehrinhalte war das entscheidend: Die Denkweise in Mengen
konnte erhalten bleiben, obwohl die technische Grundlage der Datenstruktur neu
aufgebaut werden musste.

Zugleich bereiten diese Operationen bereits die späteren, fachlich noch
wichtigeren Konstruktionen vor. Sobald Mengen zuverlässig vereinigt,
geschnitten oder voneinander abgezogen werden können, liegt der nächste Schritt
nahe: die Arbeit mit kartesischen Produkten, Potenzmengen und verschachtelten
Mengenstrukturen. Genau diesen Übergang behandelt der folgende Abschnitt.

\subsection{Kartesisches Produkt als direkte Mengenkonstruktion}
\label{subsec:rs-api-kartesisches-produkt}

Nachdem die klassischen Mengenoperationen betrachtet wurden, folgt nun eine
weitere Konstruktion, die in den Lehrnotebooks regelmäßig benötigt wird: das
kartesische Produkt. Für die theoretische Informatik ist diese Operation
besonders wichtig, weil sie auf natürliche Weise Mengen von Paaren erzeugt.
Solche Paare treten etwa bei Relationen, Übergängen oder Zustandskombinationen
auf.

Gerade hier zeigt sich erneut die Eleganz der ursprünglichen Python-Lösung.
Das kartesische Produkt wird in den Notebooks nicht über eine zusätzliche
Bibliothek konstruiert, sondern direkt als Mengenkomprehension formuliert. Das
passt sehr gut zur mathematischen Idee, weil die Zielmenge unmittelbar über
alle Paare beschrieben wird, die aus den beiden Ausgangsmengen gebildet werden
können. Python unterstützt genau diese Form der Set-Comprehension als direkte
Konstruktion einer Menge aus Iterationen und optionalen Bedingungen
\autocite{python_expressions,python_set_comprehension}.

Ein typisches Beispiel aus den Notebooks lautet:

\begin{listing}[H]
	\caption{Kartesisches Produkt in Python als Mengenkomprehension}
	\label{lst:python_cart_prod}
	\begin{minted}{python}
		def cart_prod[S, T](A: set[S], B: set[T]) -> set[tuple[S, T]]:
		return { (x, y) for x in A for y in B }
	\end{minted}
\end{listing}

Didaktisch ist diese Schreibweise besonders gelungen. Die mathematische
Definition des kartesischen Produkts wird fast direkt in Code übersetzt:
Für jedes \(x\) aus \(A\) und jedes \(y\) aus \(B\) wird das Paar \((x, y)\)
gebildet, und alle diese Paare werden zu einer Menge zusammengefasst. Gerade
deshalb war es für die Portierung wichtig, auch in TypeScript wieder eine
Mengenoperation bereitzustellen, die genau diesen fachlichen Schritt direkt
ausdrückt.

Mit \mintinline{text}{RecursiveSet} wird das kartesische Produkt deshalb nicht
über verschachtelte Schleifen in Arrays oder über Hilfskonstruktionen
formuliert, sondern als eigene Mengenoperation angeboten:

\begin{listing}[H]
	\caption{Kartesisches Produkt mit \texttt{RecursiveSet}}
	\label{lst:ts_cartesian_product}
	\begin{minted}{typescript}
		const A = new RecursiveSet(1, 2);
		const B = new RecursiveSet("a", "b");
		
		const result = A.cartesianProduct(B);
		console.log(result); // {(1, a), (1, b), (2, a), (2, b)}
	\end{minted}
\end{listing}

Die Oberflächensyntax unterscheidet sich hier sichtbar von Python. Während
Python die Konstruktion über eine Mengenkomprehension beschreibt, verwendet
\mintinline{text}{RecursiveSet} die explizite Methode
\mintinline{text}{cartesianProduct}. Fachlich bleibt der Ausdruck jedoch
derselbe: Das Ergebnis ist eine Menge aller geordneten Paare aus den beiden
Ausgangsmengen.

Zugleich zeigt dieses Beispiel sehr deutlich, warum \mintinline{text}{Tuple}
in der Bibliothek benötigt wird. Das Ergebnis des kartesischen Produkts ist
nicht einfach eine lose Sammlung beliebiger Zweiergruppen, sondern eine Menge
strukturell vergleichbarer Paare. Genau dafür dient \mintinline{text}{Tuple}
als stabiler, hashbarer und strukturell vergleichbarer Baustein.

Für die Portierung der Notebooks ist das ein wichtiger Gewinn. Die Python-
Version formuliert das kartesische Produkt direkt auf Mengenebene. Mit
\mintinline{text}{RecursiveSet} kann dieselbe fachliche Idee auch in
TypeScript wieder unmittelbar auf Mengenebene ausgedrückt werden. Damit bleibt
die mathematische Denkweise des Lehrmaterials erhalten, obwohl die technische
Umsetzung im Hintergrund vollständig neu aufgebaut werden musste.

Gerade an diesem Beispiel wird außerdem sichtbar, dass
\mintinline{text}{RecursiveSet} nicht nur einfache Mengenoperationen wie
Vereinigung oder Differenz unterstützt, sondern auch neue Mengen aus
zusammengesetzten Werten erzeugen kann. Das ist für die weiteren Kapitel
wichtig, weil viele Algorithmen der theoretischen Informatik nicht nur mit
einzelnen Werten, sondern mit Mengen von Paaren und anderen strukturierten
Objekten arbeiten.


\subsection{Transformationen auf Mengen}
\label{subsec:rs-api-transformationen}

Die bisher betrachteten Operationen wie Vereinigung, Schnitt oder kartesisches
Produkt sind unmittelbare mathematische Mengenkonstruktionen. Daneben bietet
\mintinline{text}{RecursiveSet} jedoch noch eine zweite, eher funktionale
Arbeitsweise an. Diese war in den ursprünglichen Python-Notebooks nicht der
zentrale Stil, eröffnet in der TypeScript-Portierung aber zusätzliche
Möglichkeiten, Mengen elegant und prägnant weiterzuverarbeiten.

Didaktisch ist dieser Abschnitt deshalb wichtig, weil hier nicht mehr nur
gefragt wird, \enquote{Welche Menge entsteht aus zwei anderen Mengen?}, sondern
auch: \enquote{Wie lässt sich eine vorhandene Menge systematisch
	transformieren, filtern oder auswerten?} Gerade bei Algorithmen, die viele
Zwischenschritte auf Mengen ausführen, kann diese Arbeitsweise den Code
deutlich lesbarer machen.

Im Unterschied zu den bisherigen Beispielen steht der Python-Vergleich hier
etwas weniger im Vordergrund. In Python wurden solche Schritte in den
Notebooks eher über Schleifen oder Mengenkomprehensionen formuliert. Mit
\mintinline{text}{RecursiveSet} werden dafür nun bewusst eigene Methoden
bereitgestellt. Die Datenstruktur nähert sich damit in ihrer Benutzung an
funktionale Collection-APIs an, ohne ihre eigentliche Rolle als Menge zu
verlieren.

\paragraph*{\texttt{map}: Werte einer Menge abbilden}
Mit \mintinline{text}{map} wird jedes Element einer Menge durch eine Funktion
auf einen neuen Wert abgebildet. Das Ergebnis ist wiederum eine Menge. Diese
Operation ist besonders dann nützlich, wenn aus einer Menge von Werten eine
Menge strukturell verwandter Werte erzeugt werden soll.

\begin{listing}[H]
	\caption{Abbildung einer Menge mit \texttt{map}}
	\label{lst:rs-map}
	\begin{minted}{typescript}
		const states = new RecursiveSet("q0", "q1", "q2");
		
		const labeled = states.map(state => `State:${state}`);
		console.log(labeled); // {State:q0, State:q1, State:q2}
	\end{minted}
\end{listing}

Wichtig ist hier, dass das Ergebnis wieder eine Menge ist. Dadurch bleiben die
Eigenschaften einer Mengenstruktur erhalten. Entstehen bei der Abbildung
mehrfach dieselben Zielwerte, so werden diese nicht mehrfach gespeichert,
sondern wie gewohnt zusammengefasst.

\paragraph*{\texttt{filter}: eine Teilmenge auswählen}
Mit \mintinline{text}{filter} wird aus einer vorhandenen Menge genau die
Teilmenge ausgewählt, deren Elemente eine Bedingung erfüllen. Fachlich ist das
eng verwandt mit der mathematischen Mengenschreibweise
\enquote{alle \(x \in M\) mit Eigenschaft \(P(x)\)}.

\begin{listing}[H]
\caption{Auswahl einer Teilmenge mit \texttt{filter}}
\label{lst:rs-filter}
\begin{minted}{typescript}
	const numbers = new RecursiveSet(1, 2, 3, 4, 5, 6);
	
	const even = numbers.filter(n => n % 2 === 0);
	console.log(even); // {2, 4, 6}
\end{minted}
\end{listing}

Gerade diese Operation passt sehr gut zur Denkweise der theoretischen
Informatik. Häufig interessiert nicht die gesamte Menge, sondern nur die
Teilmenge erreichbarer, akzeptierender oder anderweitig relevanter Elemente.
\mintinline{text}{filter} macht genau diesen Schritt unmittelbar sichtbar.

\paragraph*{\texttt{reduce}: eine Menge zu einem einzelnen Wert zusammenfassen}
Während \mintinline{text}{map} und \mintinline{text}{filter} wieder Mengen
erzeugen, führt \mintinline{text}{reduce} aus der Mengenwelt heraus. Alle
Elemente werden schrittweise zu einem einzelnen Ergebniswert zusammengeführt.
Das ist besonders nützlich, wenn aus einer Menge etwa eine Zählung, eine
Zeichenkette oder eine aggregierte Kennzahl gebildet werden soll.

\begin{listing}[H]
\caption{Zusammenfassung einer Menge mit \texttt{reduce}}
\label{lst:rs-reduce}
\begin{minted}{typescript}
const numbers = new RecursiveSet(1, 2, 3, 4);

const sum = numbers.reduce((acc, n) => acc + n, 0);
console.log(sum); // 10
\end{minted}
\end{listing}

Für die Portierung ist diese Methode vor allem deshalb interessant, weil sie
eine strukturierte Alternative zu handgeschriebenen Schleifen bietet. Wenn
eine Menge nur noch ausgewertet und nicht mehr als Menge weiterverwendet wird,
ist \mintinline{text}{reduce} oft die prägnanteste Formulierung.

\paragraph*{\texttt{every} und \texttt{some}: Eigenschaften ganzer Mengen prüfen}
Mit \mintinline{text}{every} und \mintinline{text}{some} lässt sich prüfen, ob
alle beziehungsweise mindestens ein Element einer Menge eine Bedingung
erfüllen. Diese beiden Methoden sind besonders dann praktisch, wenn eine
Mengenstruktur nicht transformiert, sondern nur logisch bewertet werden soll.

\begin{listing}[H]
\caption{Logische Prüfungen mit \texttt{every} und \texttt{some}}
\label{lst:rs-every-some}
\begin{minted}{typescript}
const numbers = new RecursiveSet(2, 4, 6, 7);

console.log(numbers.every(n => n % 2 === 0)); // false
console.log(numbers.some(n => n % 2 !== 0));  // true
\end{minted}
\end{listing}

Solche Prüfungen wirken auf den ersten Blick unscheinbar, sind im praktischen
Code aber sehr hilfreich. Statt eine Schleife mit manueller Abbruchlogik zu
schreiben, kann direkt ausgedrückt werden, welche Eigenschaft die Menge als
Ganzes erfüllen soll.

\paragraph*{Warum diese Methoden für die Portierung nützlich sind}
Gerade dieser funktionale Stil war nicht der dominierende Ausgangspunkt der
ursprünglichen Python-Notebooks. Dennoch erweitert er die Portierung in einer
sinnvollen Weise. Viele wiederkehrende Verarbeitungsschritte auf Mengen lassen
sich dadurch kompakter und klarer formulieren als mit expliziten Schleifen.

Für die Studienarbeit ist das ein wichtiger Unterschied zu den bisherigen
Abschnitten. Bei Vereinigung, Schnitt oder kartesischem Produkt stand vor
allem die mathematische Mengenlehre im Vordergrund. Bei
\mintinline{text}{map}, \mintinline{text}{filter} oder \mintinline{text}{reduce}
kommt nun zusätzlich ein moderner Collection-Stil hinzu, der in TypeScript
besonders gut zur restlichen Umgebung passt.

Gleichzeitig bleibt die zentrale Eigenschaft von \mintinline{text}{RecursiveSet}
erhalten: Auch diese Methoden arbeiten nicht mit beliebigen Listen, sondern
mit Mengen. Die Ergebnisse bleiben also weiterhin duplikatfrei und für
zusammengesetzte Werte strukturell vergleichbar. Genau dadurch verbinden diese
Transformationen funktionale Eleganz mit der fachlichen Präzision, die für die
Portierung der Lehrinhalte benötigt wird.

Im nächsten Abschnitt wird diese Linie noch weitergeführt. Dort geht es mit
\mintinline{text}{filterMap} und \mintinline{text}{flatMap} um zwei Methoden,
die mehrere Verarbeitungsschritte bewusst zusammenfassen und gerade für
größere Algorithmen besonders interessant sind.


\subsection{Effiziente Kombinationsschritte mit \texttt{filterMap} und \texttt{flatMap}}
\label{subsec:rs-api-filtermap-flatmap}

Die bisher betrachteten Methoden \mintinline{text}{map}, \mintinline{text}{filter}
und \mintinline{text}{reduce} decken bereits viele typische Verarbeitungsschritte
ab. In der Praxis treten jedoch häufig zwei wiederkehrende Muster auf. Zum einen
sollen nur bestimmte Elemente berücksichtigt und diese anschließend direkt in
eine neue Form überführt werden. Zum anderen wird für viele Eingabewerte jeweils
eine neue Teilmenge erzeugt, deren Inhalte am Ende in einer gemeinsamen Zielmenge
landen sollen.

Genau für diese beiden Muster stellt \mintinline{text}{RecursiveSet} die
Methoden \mintinline{text}{filterMap} und \mintinline{text}{flatMap} bereit.
Beide Methoden sind nicht nur bequeme Kurzschreibweisen, sondern fassen
häufige Kombinationen von Verarbeitungsschritten direkt zu einer eigenen
Operation zusammen.

Didaktisch ist dabei wichtig, dass dieser Abschnitt stärker von den bisherigen
Python-Vergleichen abweicht. In Python werden vergleichbare Schritte häufig
über Schleifen oder Mengenkomprehensionen formuliert. Python kennt gerade für
Mengen eine sehr kompakte Schreibweise der Form
\mintinline{python}{{expression for x in iterable if condition}}
\autocite{python_expressions,python_set_comprehension}. Mit
\mintinline{text}{RecursiveSet} wird dieselbe Grundidee nun zusätzlich in
Methodenform angeboten.

\paragraph*{\texttt{filterMap}: Auswählen und direkt umformen}
Die Methode \mintinline{text}{filterMap} verbindet zwei gedanklich aufeinander
folgende Schritte: Zuerst wird geprüft, welche Elemente einer Menge überhaupt
relevant sind, danach werden genau diese Elemente in eine neue Form
überführt.

Ein einfaches Beispiel zeigt das Prinzip. Aus einer Menge von Zahlen sollen
nur die geraden Werte übernommen und direkt in beschreibende Strings
umgewandelt werden:

\begin{listing}[H]
	\caption{Filtern und Abbilden mit \texttt{filterMap}}
	\label{lst:rs-filtermap-generic}
	\begin{minted}{typescript}
		const numbers = new RecursiveSet(1, 2, 3, 4, 5, 6);
		
		const labels = numbers.filterMap(
		n => n % 2 === 0,
		n => `even:${n}`
		);
		
		console.log(labels); // {even:2, even:4, even:6}
	\end{minted}
\end{listing}

Natürlich ließe sich dieselbe Idee auch als
\mintinline{text}{numbers.filter(...).map(...)} schreiben. Der Vorteil von
\mintinline{text}{filterMap} besteht jedoch darin, dass kein zusätzliches
Zwischenergebnis aufgebaut werden muss. Die Auswahl und die Umformung werden
bewusst zu einem einzigen Verarbeitungsschritt zusammengezogen.

\paragraph*{\texttt{flatMap}: Viele Teilmengen direkt zusammenführen}
Noch interessanter ist \mintinline{text}{flatMap}. Die Grundidee ist hier:
Für jedes Element einer Eingabe wird eine neue Menge erzeugt, und alle diese
Teilmengen werden anschließend direkt in eine gemeinsame Zielmenge übernommen.
Diese Idee ist eng verwandt mit der bekannten JavaScript-Intuition von
\mintinline{text}{flatMap}, die sich als \mintinline{text}{map(...).flat()}
mit einer Ebene des Zusammenführens verstehen lässt \autocite{mdn_array_flatmap}.

Ein generisches Beispiel dafür ist die Erzeugung aller Präfixe einer Menge von
Wörtern. Für jedes Wort entsteht zunächst eine eigene Menge seiner Präfixe.
Diese Teilmengen sollen danach zu einer einzigen Ergebnismenge vereinigt
werden.

Zunächst kann man das noch mit wiederholten Vereinigungen formulieren:

\begin{listing}[H]
	\caption{Wiederholte Vereinigungen als Ausgangspunkt}
	\label{lst:rs-flatmap-union-generic}
	\begin{minted}{typescript}
		function prefixes(word: string): RecursiveSet<string> {
			const result = new RecursiveSet<string>();
			
			for (let i = 1; i <= word.length; i++) {
				result.add(word.slice(0, i));
			}
			
			return result;
		}
		
		function allPrefixes(words: string[]): RecursiveSet<string> {
			let result = new RecursiveSet<string>();
			
			for (const word of words) {
				result = result.union(prefixes(word));
			}
			
			return result;
		}
		
		console.log(allPrefixes(["cat", "car"]));
		// {c, ca, cat, car}
	\end{minted}
\end{listing}

Diese Variante ist leicht verständlich und fachlich korrekt. Gleichzeitig
zeigt sie aber bereits das wiederkehrende Muster: Für jedes Eingabewort wird
eine Teilmenge erzeugt, und danach wird diese Teilmenge mit der bisherigen
Gesamtmenge vereinigt.

Genau dieses Muster lässt sich mit \mintinline{text}{flatMap} direkter
ausdrücken:

\begin{listing}[H]
	\caption{Direktes Einsammeln mit \texttt{flatMap}}
	\label{lst:rs-flatmap-generic}
	\begin{minted}{typescript}
		function prefixes(word: string): RecursiveSet<string> {
			const result = new RecursiveSet<string>();
			
			for (let i = 1; i <= word.length; i++) {
				result.add(word.slice(0, i));
			}
			
			return result;
		}
		
		function allPrefixes(words: string[]): RecursiveSet<string> {
			const result = new RecursiveSet<string>();
			result.flatMap(words, prefixes);
			return result;
		}
		
		console.log(allPrefixes(["cat", "car"]));
		// {c, ca, cat, car}
	\end{minted}
\end{listing}

Der Vorteil dieser Formulierung liegt nicht nur in kürzerem Code. Viel
wichtiger ist, dass die eigentliche Struktur des Verfahrens nun klarer
sichtbar wird: Es gibt eine Quelle von Eingabewerten, eine Funktion, die aus
jedem Eingabewert eine Teilmenge erzeugt, und eine Zielmenge, in die alle
Ergebnisse direkt eingesammelt werden.

Gerade deshalb ist \mintinline{text}{flatMap} in \mintinline{text}{RecursiveSet}
mutierend implementiert. Die Methode arbeitet auf einer bereits vorhandenen
Zielmenge und erweitert diese Schritt für Schritt. Das unterscheidet sie
bewusst von rein funktionalen Varianten, die nach jedem Verarbeitungsschritt
wieder eine neue Gesamtstruktur erzeugen würden.

Für die Portierung ist das ein sinnvoller Mittelweg. Einerseits bleibt die
Idee sehr nah an einem funktionalen Transformationsstil. Andererseits wird ein
häufiger Aufbau unnötiger Zwischenmengen vermieden. Genau dadurch eignet sich
\mintinline{text}{flatMap} besonders gut für größere Konstruktionsschritte, bei
denen viele kleine Ergebnismengen zu einer einzigen Gesamtmenge zusammengeführt
werden müssen.

Zusammen zeigen \mintinline{text}{filterMap} und \mintinline{text}{flatMap},
dass \mintinline{text}{RecursiveSet} nicht nur das vertraute Mengenverhalten
aus Python wiederherstellt, sondern darüber hinaus auch eine moderne und für
TypeScript gut passende Arbeitsoberfläche anbietet. Die mathematische Idee der
Menge bleibt erhalten, wird aber um prägnante Verarbeitungsschritte ergänzt,
die den Code an vielen Stellen klarer und direkter machen.


\subsection{Zwischen \texttt{set} und \texttt{frozenset}: Der Lebenszyklus von \texttt{RecursiveSet}}
\label{subsec:rs-api-freeze-frozenset}

Die bisher betrachteten Beispiele haben \mintinline{text}{RecursiveSet}
zunächst wie eine gewöhnliche veränderliche Menge erscheinen lassen. Elemente
konnten mit \mintinline{text}{add} eingefügt und mit \mintinline{text}{remove}
wieder entfernt werden. Für viele Nutzungssituationen ist genau dieses
Verhalten auch erwünscht. Gleichzeitig benötigt die Bibliothek jedoch an
bestimmten Stellen eine stabile, unveränderliche Mengenstruktur, etwa wenn
eine Menge selbst wieder als Wert in einer anderen Menge gespeichert wird.

Gerade hier lohnt sich der Vergleich mit Python. Python unterscheidet zwischen
\mintinline{python}{set} und \mintinline{python}{frozenset}. Während
\mintinline{python}{set} veränderlich ist, ist \mintinline{python}{frozenset}
unveränderlich und hashbar \autocite{python_frozenset_builtin,python_set_objects}.
Genau deshalb kann ein \mintinline{python}{frozenset} als Element einer
anderen Menge oder als Schlüssel in einem Dictionary verwendet werden
\autocite{python_set_objects,python_set_types_legacy}.

Ein kleines Python-Beispiel zeigt diesen Unterschied unmittelbar:

\begin{listing}[H]
	\caption{Python unterscheidet zwischen \texttt{set} und \texttt{frozenset}}
	\label{lst:python_set_frozenset}
	\begin{minted}{python}
		A = {1, 2}
		B = frozenset({1, 2})
		
		A.add(3)
		print(A)  # {1, 2, 3}
		
		# B.add(3) wäre nicht erlaubt
	\end{minted}
\end{listing}

Didaktisch ist diese Trennung in Python sehr klar. Schon am verwendeten Typ
ist erkennbar, ob eine Menge später noch verändert werden darf oder nicht.
Gerade für verschachtelte Mengen ist \mintinline{python}{frozenset} deshalb
eine naheliegende Lösung.

Für die TypeScript-Portierung wäre eine direkte Nachbildung mit zwei
vollständig getrennten Konstrukten jedoch unhandlich gewesen. Die Benutzung
der Bibliothek sollte sich möglichst einheitlich anfühlen und nicht ständig
zwischen einer \enquote{aufbaubaren} und einer \enquote{fertigen} Mengenform
wechseln. Deshalb verfolgt \mintinline{text}{RecursiveSet} einen anderen
Ansatz: Statt zwei verschiedene Typen bereitzustellen, arbeitet die
Bibliothek mit einem Lebenszyklus derselben Datenstruktur.

Eine \mintinline{text}{RecursiveSet}-Instanz ist zunächst veränderlich. Sie
kann normal aufgebaut, erweitert und reduziert werden. Erst dann, wenn ihre
Struktur stabil sein muss, wird sie eingefroren. Genau dieser Übergang macht
aus einer zunächst bearbeitbaren Menge eine unveränderliche Struktur, die
sicher in anderen rekursiven Strukturen verwendet werden kann.

Das folgende Beispiel zeigt dieses Verhalten:

\begin{listing}[H]
	\caption{\texttt{RecursiveSet} verbindet beide Rollen über Einfrieren}
	\label{lst:ts_recursive_set_freeze}
	\begin{minted}{typescript}
		const A = new RecursiveSet(1, 2);
		A.add(3);
		
		const B = new RecursiveSet(A);
		
		// Ab hier ist A eingefroren und darf nicht mehr verändert werden.
		// A.add(4); // Fehler: Frozen Set modified.
	\end{minted}
\end{listing}

An diesem Beispiel ist die fachliche Idee besonders wichtig. Die Menge
\mintinline{text}{A} beginnt als gewöhnliche veränderliche Menge. Sobald sie
jedoch selbst als Element einer anderen Menge verwendet wird, muss ihre
Struktur stabil bleiben. Genau an dieser Stelle wird sie eingefroren. Damit
übernimmt \mintinline{text}{RecursiveSet} gewissermaßen die Rolle, die in
Python von \mintinline{python}{frozenset} erfüllt wird, ohne dafür einen
eigenen zweiten Nutzertyp einzuführen.

Diese Designentscheidung war für die Portierung besonders vorteilhaft.
Einerseits bleibt die API einfach, weil überall dieselbe Datenstruktur
verwendet wird. Andererseits geht die fachlich wichtige Unveränderlichkeit
nicht verloren. Die Bibliothek verbindet also die Flexibilität einer
aufbaubaren Menge mit der Stabilität einer unveränderlichen Mengenstruktur.

Falls eine eingefrorene Menge später doch noch verändert werden soll, stellt
die Bibliothek mit \mintinline{text}{mutableCopy()} einen bewussten Rückweg
bereit. Anstatt ein einmal eingefrorenes Objekt wieder direkt veränderlich zu
machen, wird eine neue, bearbeitbare Kopie erzeugt:

\begin{listing}[H]
	\caption{Weiterarbeiten mit einer veränderlichen Kopie}
	\label{lst:ts_recursive_set_mutable_copy}
	\begin{minted}{typescript}
		const A = new RecursiveSet(1, 2);
		const B = new RecursiveSet(A);
		
		// A ist nun eingefroren
		const C = A.mutableCopy();
		C.add(3);
		
		console.log(C); // {1, 2, 3}
	\end{minted}
\end{listing}

Gerade dieser Mechanismus zeigt, dass das Einfrieren nicht als lästige
Einschränkung verstanden werden sollte. Es handelt sich vielmehr um eine
bewusste Sicherheitsregel, die stabil macht, was stabil sein muss, ohne den
praktischen Arbeitsfluss unnötig zu erschweren. Wer weiter an einer Menge
arbeiten möchte, erhält mit \mintinline{text}{mutableCopy()} eine klare und
kontrollierte Möglichkeit dazu.

Damit verbindet \mintinline{text}{RecursiveSet} zwei Anforderungen, die in
Python auf zwei getrennte Konstrukte verteilt sind. Die Menge kann zunächst
wie ein \mintinline{python}{set} aufgebaut werden, übernimmt bei Bedarf aber
die Rolle eines \mintinline{python}{frozenset}. Für die Portierung erwies
sich genau dieser Lebenszyklus als die elegantere Lösung, weil er die
Benutzung vereinheitlicht und zugleich die für rekursive Mengenstrukturen
notwendige Stabilität sichert.

Im nächsten Abschnitt wird dieser API-Blick noch einmal von einer anderen
Seite ergänzt. Dort geht es nicht mehr um den Lebenszyklus der Menge, sondern
um die Frage, welche Arten von Werten überhaupt in \mintinline{text}{RecursiveSet}
gespeichert werden dürfen und welche Verträge diese Werte erfüllen müssen.


\subsection{Welche Werte überhaupt gespeichert werden dürfen}
\label{subsec:rs-api-wertetypen}

Bisher wurde \mintinline{text}{RecursiveSet} vor allem als Mengenstruktur mit
fachlich passendem Verhalten betrachtet. Für die praktische Nutzung stellt sich
nun jedoch eine naheliegende Frage: Welche Arten von Werten dürfen überhaupt in
einer solchen Menge gespeichert werden?

Diese Frage ist für unsere Bibliothek besonders wichtig, weil
\mintinline{text}{RecursiveSet} nicht einfach beliebige JavaScript-Werte
akzeptiert. Die Datenstruktur soll nicht nur irgendetwas speichern, sondern
Werte, deren Gleichheit und Hash-Verhalten zuverlässig bestimmt werden können.
Gerade deshalb arbeitet die Bibliothek mit einem klaren Wertevertrag.

Im Kern unterscheidet die API zwei Gruppen erlaubter Werte. Zum einen dürfen
primitive Werte gespeichert werden, in unserem Fall also Zahlen und Strings.
Zum anderen dürfen komplexere Werte gespeichert werden, sofern sie den
Vertrag \mintinline{text}{Structural} erfüllen. Dieser Vertrag verlangt, dass
ein Wert strukturelle Gleichheit, einen stabilen Hash-Wert und eine
deterministische Textdarstellung bereitstellt.

Die zugehörigen Typdefinitionen lauten:

\begin{listing}[H]
	\caption{Erlaubte Wertarten in \texttt{RecursiveSet}}
	\label{lst:rs-value-types}
	\begin{minted}{typescript}
		type Primitive = string | number;
		type Value = Primitive | Structural;
		
		interface Structural {
			readonly hashCode: number;
			equals(other: unknown): boolean;
			toString(): string;
		}
	\end{minted}
\end{listing}

Didaktisch ist hier vor allem wichtig, dass nicht \enquote{komplex} gegen
\enquote{einfach} unterschieden wird, sondern \enquote{direkt verständlich}
gegen \enquote{explizit beschrieben}. Zahlen und Strings kann die Bibliothek
selbst unmittelbar behandeln. Bei komplexeren Strukturen muss dagegen klar
festgelegt sein, wie Gleichheit und Hashing funktionieren. Genau dafür dient
das Interface \mintinline{text}{Structural}. TypeScript verwendet Interfaces
gerade dazu, Verträge über die Form eines Werts zu beschreiben; entscheidend
ist also nicht der Name einer Klasse, sondern ob die geforderten Mitglieder
vorhanden sind \autocite{ts_interfaces,ts_type_compatibility}.

\paragraph*{Primitive Werte}
Am einfachsten ist die Benutzung bei primitiven Werten. Zahlen und Strings
können ohne weitere Vorbereitung direkt gespeichert werden:

\begin{listing}[H]
	\caption{Primitive Werte in \texttt{RecursiveSet}}
	\label{lst:rs-primitive-values}
	\begin{minted}{typescript}
		const numbers = new RecursiveSet(1, 2, 3);
		const states  = new RecursiveSet("q0", "q1", "q2");
		
		console.log(numbers); // {1, 2, 3}
		console.log(states);  // {q0, q1, q2}
	\end{minted}
\end{listing}

Diese Fälle wirken zunächst unspektakulär, sind aber didaktisch wichtig. Sie
zeigen, dass \mintinline{text}{RecursiveSet} nicht nur für komplizierte
Spezialstrukturen gedacht ist. Auch gewöhnliche Mengen aus Zahlen oder
Bezeichnern lassen sich direkt und ohne zusätzlichen Aufwand formulieren.

\paragraph*{Strukturelle Werte}
Interessant wird die Datenstruktur bei Werten, die selbst aus mehreren
Bestandteilen bestehen. Genau dafür genügt ein primitiver Typ nicht mehr. Ein
solcher Wert muss dann den Vertrag \mintinline{text}{Structural} erfüllen.

Ein typisches Beispiel ist \mintinline{text}{Tuple}. Das Tupel ist kein
eingebauter primitiver Wert, besitzt aber genau die benötigten Eigenschaften:
Es kann Gleichheit strukturell prüfen, einen Hash-Wert liefern und eine
stabile Darstellung erzeugen.

\begin{listing}[H]
	\caption{Ein struktureller Wert mit \texttt{Tuple}}
	\label{lst:rs-structural-tuple}
	\begin{minted}{typescript}
		const edges = new RecursiveSet<Tuple<[string, string]>>();
		edges.add(new Tuple("q0", "q1"));
		edges.add(new Tuple("q1", "q2"));
		
		console.log(edges); // {(q0, q1), (q1, q2)}
	\end{minted}
\end{listing}

Gerade an solchen Beispielen wird sichtbar, warum die Bibliothek nicht einfach
beliebige Objekte zulässt. Bei einem Tupel ist präzise definiert, wann zwei
Tupel gleich sind. Dasselbe gilt für ihren Hash-Wert und ihre Darstellung.
Damit wird aus einem zusammengesetzten Wert ein sauber beschriebener
Mengenbaustein.

\paragraph*{Rekursive Werte}
Eine besondere Stärke von \mintinline{text}{RecursiveSet} besteht darin, dass
auch rekursive Strukturen unterstützt werden. Eine Menge kann also selbst
wieder andere Mengen enthalten, sofern diese die nötigen Bedingungen
erfüllen.

\begin{listing}[H]
	\caption{Mengen von Mengen}
	\label{lst:rs-set-of-sets}
	\begin{minted}{typescript}
		const A = new RecursiveSet(1, 2);
		const B = new RecursiveSet(2, 1);
		
		const meta = new RecursiveSet<RecursiveSet<number>>();
		meta.add(A);
		meta.add(B);
		
		console.log(meta);      // {{1, 2}}
		console.log(meta.size); // 1
	\end{minted}
\end{listing}

Dieses Beispiel ist für die Arbeit besonders aufschlussreich. Die innere Menge
ist hier selbst wieder ein struktureller Wert. Dadurch kann sie in einer
äußeren Menge gespeichert und dort nach ihrem Inhalt statt nach ihrer
bloßen Identität behandelt werden.

\paragraph*{Warum Plain Objects ausgeschlossen sind}
Besonders wichtig ist nun die Frage, warum gewöhnliche JavaScript-Objekte der
Form \mintinline{text}{{ key: value }} nicht einfach ebenfalls zugelassen
werden. Auf den ersten Blick wirken solche Objekte sehr flexibel. Für unsere
Bibliothek wäre diese Offenheit jedoch gerade problematisch.

JavaScript vergleicht nicht-primitive Objekte nicht nach ihrer inhaltlichen
Struktur, sondern nach ihrer Identität. Zwei verschiedene Objekte mit
denselben Feldern gelten daher ohne zusätzliche Logik nicht automatisch als
gleich \autocite{mdn_equality_sameness}. Genau deshalb würde ein beliebiges
Plain Object noch nicht die Informationen mitbringen, die
\mintinline{text}{RecursiveSet} für Wertsemantik benötigt.

Das Problem ist also nicht, dass Plain Objects \enquote{zu komplex} wären.
Das Problem ist vielmehr, dass für sie in unserer Bibliothek nicht von selbst
klar ist, wie Gleichheit, Hashing und stabile Darstellung funktionieren
sollen. Diese Entscheidungen müssen explizit beschrieben werden. Genau
deshalb fordert die API für komplexe Werte den Vertrag
\mintinline{text}{Structural}.

Ein absichtlich nicht unterstütztes Beispiel sieht daher so aus:

\begin{listing}[H]
	\caption{Ein gewöhnliches Objekt ist kein zulässiger struktureller Wert}
	\label{lst:rs-plain-object-not-supported}
	\begin{minted}{typescript}
		// Nicht vorgesehen:
		const wrong = new RecursiveSet<{ left: number; right: number }>();
		
		// Für Plain Objects ist in der Bibliothek nicht festgelegt,
		// wie hashCode, equals und eine stabile Darstellung aussehen sollen.
	\end{minted}
\end{listing}

Für die Portierung ist genau diese Einschränkung kein Nachteil, sondern eine
bewusste Vereinfachung. Die Bibliothek akzeptiert nicht beliebige dynamische
Objektformen, sondern nur solche Werte, deren Verhalten als Mengenelement
sauber beschrieben ist. Das macht die Benutzung zwar etwas strenger, erhöht
aber die fachliche Verlässlichkeit der gesamten Datenstruktur.

\paragraph*{Einordnung für die Studienarbeit}
Gerade dieser Abschnitt zeigt noch einmal sehr deutlich den Unterschied
zwischen der bisherigen Python-Welt und der neuen TypeScript-Lösung. In den
Python-Notebooks konnten viele Datenformen sehr frei entstehen. In der
TypeScript-Portierung wird nun bewusster entschieden, welche Werte als echte
Mengenelemente geeignet sind.

Für unsere Studienarbeit ist das ein klarer Vorteil. Primitive Werte wie
Zahlen und Strings bleiben unkompliziert nutzbar. Zusammengesetzte Werte wie
Tupel, Mengen oder spätere Spezialstrukturen können ebenfalls gespeichert
werden, müssen ihr Verhalten jedoch explizit über \mintinline{text}{Structural}
beschreiben. Dadurch bleibt die Bibliothek offen für fachlich sinnvolle
Strukturen, ohne in beliebige und schwer kontrollierbare Objektformen
abzugleiten.

Im nächsten Abschnitt werden diese erlaubten Wertarten noch um eine weitere
Perspektive ergänzt. Dort geht es nicht mehr darum, \emph{welche} Werte
gespeichert werden dürfen, sondern um die Regeln und Grenzen, die für die
korrekte Benutzung der Bibliothek insgesamt gelten.


\subsection{Regeln und Grenzen der Bibliothek}
\label{subsec:rs-api-regeln}

Die bisherigen Abschnitte haben \mintinline{text}{RecursiveSet} vor allem als
leistungsfähige und fachlich passende Mengenstruktur eingeführt. Damit diese
Eigenschaften in der Praxis tatsächlich erhalten bleiben, arbeitet die
Bibliothek jedoch nicht völlig voraussetzungslos. Bestimmte Nutzungsregeln
müssen eingehalten werden, damit Gleichheit, Hashing und verschachtelte
Strukturen zuverlässig funktionieren.

Didaktisch ist an dieser Stelle ein wichtiger Perspektivwechsel nötig. Die
folgenden Punkte sind nicht als bloße Verbotsliste zu verstehen. Sie machen
vielmehr sichtbar, welche Annahmen die Datenstruktur über ihre Elemente trifft.
Gerade weil \mintinline{text}{RecursiveSet} nicht nur primitive Werte, sondern
auch zusammengesetzte und rekursive Strukturen unterstützt, müssen einige
Grundbedingungen klar festgelegt werden.

\paragraph*{Nur endliche Zahlen}
Die Bibliothek erlaubt als primitive numerische Werte nur endliche Zahlen.
Werte wie \mintinline{text}{NaN}, \mintinline{text}{Infinity} oder
\mintinline{text}{-Infinity} sind bewusst ausgeschlossen. JavaScript behandelt
\mintinline{text}{NaN} in Vergleichen besonders, denn \mintinline{text}{NaN}
ist nicht einmal zu sich selbst gleich, und \mintinline{text}{Number.isFinite}
liefert für \mintinline{text}{NaN} sowie für positive und negative
Unendlichkeit jeweils \mintinline{text}{false} \autocite{mdn_nan,mdn_number_isfinite}.

Für die Bibliothek wäre dies problematisch, weil numerische Werte in Hashing
und Gleichheitsprüfungen möglichst stabil und eindeutig behandelbar sein
müssen. Endliche Zahlen erfüllen diese Erwartung. Spezialwerte wie
\mintinline{text}{NaN} oder Unendlichkeiten würden dagegen Sonderfälle in die
Mengenlogik hineintragen, die dem Ziel einer klaren und performanten
Wertsemantik widersprechen.

\begin{listing}[H]
	\caption{Zulässige und unzulässige numerische Werte}
	\label{lst:rs-finite-numbers-only}
	\begin{minted}{typescript}
		const ok = new RecursiveSet(1, 2, 3);
		
		// Nicht vorgesehen:
		// const wrong = new RecursiveSet(NaN, Infinity);
	\end{minted}
\end{listing}

\paragraph*{Keine beliebigen Plain Objects}
Wie bereits im vorigen Abschnitt vorbereitet, akzeptiert die Bibliothek
komplexe Werte nicht in beliebiger Form. Gewöhnliche JavaScript-Objekte der
Form \mintinline{text}{{ key: value }} sind nicht automatisch geeignete
Mengenelemente. JavaScript vergleicht verschiedene nicht-primitive Objekte
auch dann nicht als gleich, wenn sie dieselbe Struktur besitzen, solange sie
nicht dieselbe Objektidentität haben \autocite{mdn_equality_sameness}.

Für \mintinline{text}{RecursiveSet} reicht eine solche Referenzgleichheit
gerade nicht aus. Die Bibliothek benötigt für komplexe Werte ein explizit
beschriebenes Verhalten für Gleichheit, Hashing und Darstellung. Genau
deshalb müssen zusammengesetzte Werte den Vertrag
\mintinline{text}{Structural} erfüllen.

\begin{listing}[H]
	\caption{Plain Objects sind nicht als Mengenelemente vorgesehen}
	\label{lst:rs-no-plain-objects}
	\begin{minted}{typescript}
		// Nicht vorgesehen:
		const wrong = new RecursiveSet<{ left: number; right: number }>();
		
		// Stattdessen sollte ein eigener struktureller Typ verwendet werden.
	\end{minted}
\end{listing}

\paragraph*{Die Qualität von \texttt{hashCode} ist wichtig}
Für strukturelle Werte genügt es nicht, irgendeinen Hash-Wert bereitzustellen.
Der Hash sollte die Werte möglichst gut verteilen. Wenn sehr viele
verschiedene Elemente denselben \mintinline{text}{hashCode} liefern, dann
landen sie intern immer wieder in denselben Bereichen der Hashtabelle. Die
Bibliothek bleibt damit zwar fachlich korrekt, verliert jedoch einen Teil
ihrer erwarteten Leistungsfähigkeit.

Didaktisch ist hier vor allem wichtig, Hashing nicht als magische Zahl zu
missverstehen. Der Hash-Wert dient nicht nur als irgendein Zusatzmerkmal,
sondern als Wegweiser für Einfügen, Suchen und Löschen. Eine schlechte
Verteilung verschlechtert deshalb nicht die inhaltliche Korrektheit, aber die
Effizienz des gesamten Verfahrens.

\begin{listing}[H]
	\caption{Ein schlechter Hash-Wert ist fachlich möglich, aber ungünstig}
	\label{lst:rs-bad-hashcode}
	\begin{minted}{typescript}
		class BadValue implements Structural {
			get hashCode(): number { return 42; }
			
			equals(other: unknown): boolean {
				return other instanceof BadValue;
			}
			
			toString(): string {
				return "BadValue";
			}
		}
	\end{minted}
\end{listing}

\paragraph*{Die Darstellung muss deterministisch bleiben}
Neben Gleichheit und Hashing spielt auch die Textdarstellung eine wichtige
Rolle. \mintinline{text}{RecursiveSet} stellt Mengen bewusst so dar, dass sie
in Ausgaben und Fehlersituationen lesbar bleiben. Damit diese Darstellung bei
verschachtelten Strukturen stabil ist, müssen benutzerdefinierte
\mintinline{text}{toString()}-Implementierungen dieselbe visuelle
Vergleichslogik respektieren wie die Bibliothek selbst.

Gerade dieser Punkt ist leicht zu unterschätzen. Für die eigentliche
Mengenlogik ist die Reihenfolge von Elementen nicht entscheidend, denn Mengen
sind fachlich ungeordnet. Für die Darstellung ist eine stabile Reihenfolge
jedoch äußerst hilfreich, weil gleiche Strukturen dann auch gleich aussehen.
Die Regel dient also nicht der mathematischen Bedeutung der Menge, sondern der
Nachvollziehbarkeit ihrer Ausgabe.

\paragraph*{Eingefügte Werte dürfen sich nicht nachträglich verändern}
Sobald ein struktureller Wert in einer Menge gespeichert wurde, darf sich sein
\mintinline{text}{hashCode} nicht mehr ändern. Andernfalls würde die
Bibliothek einen Wert unter einer alten Hash-Struktur verwalten, obwohl seine
neue inhaltliche Form eigentlich an einer anderen Stelle gesucht werden
müsste. Genau deshalb ist Stabilität nach dem Einfügen keine bloße Empfehlung,
sondern eine zentrale Voraussetzung für korrektes Verhalten.

An dieser Stelle knüpft die Bibliothek an den vorigen Abschnitt zum
Einfrieren an. Dort wurde bereits aus Anwendersicht erklärt, warum eine Menge
später unveränderlich werden kann. Hier wird nun sichtbar, warum diese Regel
fachlich notwendig ist: Hash-basierte Strukturen setzen voraus, dass einmal
einsortierte Werte ihre Einordnung nicht heimlich ändern.

\paragraph*{Keine Zyklen in rekursiven Strukturen}
\mintinline{text}{RecursiveSet} unterstützt rekursive Strukturen, etwa Mengen
von Mengen. Diese Unterstützung hat jedoch eine wichtige Grenze: Zyklen sind
nicht vorgesehen. Eine Menge darf also nicht direkt oder indirekt sich selbst
enthalten.

Der Grund dafür liegt nicht in einer mathematischen Spitzfindigkeit, sondern
in der praktischen Verarbeitung. Hashing und Darstellung laufen rekursiv durch
verschachtelte Strukturen. Würde eine Menge schließlich wieder auf sich selbst
verweisen, gäbe es keinen natürlichen Endpunkt mehr für diese Verarbeitung.
Die Bibliothek prüft diesen Fall aus Leistungsgründen nicht bei jeder
Operation, sondern setzt voraus, dass solche Zyklen gar nicht erst erzeugt
werden.

\begin{listing}[H]
	\caption{Rekursive Mengen sind erlaubt, zyklische Mengen nicht}
	\label{lst:rs-no-cycles}
	\begin{minted}{typescript}
		const A = new RecursiveSet<number>(1, 2);
		const B = new RecursiveSet(A);
		
		// Nicht vorgesehen wäre eine Struktur, in der A direkt oder indirekt
		// wieder sich selbst enthält.
	\end{minted}
\end{listing}

\paragraph*{Einordnung der Regeln}
Zusammengenommen zeigen diese Regeln sehr deutlich, dass
\mintinline{text}{RecursiveSet} eine bewusst entworfene Fachstruktur ist und
kein beliebiger Container für alle möglichen JavaScript-Werte. Gerade diese
Strenge ist jedoch ein Vorteil für die Studienarbeit. Die Bibliothek macht
klar sichtbar, welche Voraussetzungen erfüllt sein müssen, damit Mengen mit
Wertsemantik zuverlässig und performant funktionieren.

Für die Portierung der Lehrnotebooks ist das ein wichtiger Gewinn. Die Regeln
engen die Bibliothek nicht willkürlich ein, sondern schaffen einen Rahmen, in
dem Gleichheit, Hashing und rekursive Mengenkonstruktionen tatsächlich
zusammenpassen. Dadurch bleibt das fachliche Verhalten der Mengen nicht nur
beobachtbar korrekt, sondern auch langfristig belastbar.

Im nächsten Abschnitt wird diese API-Perspektive noch einmal auf das gesamte
Projekt zurückgeführt. Dort geht es um die Frage, welche Rolle
\mintinline{text}{RecursiveSet} innerhalb der Portierung insgesamt spielt und
warum gerade diese Datenstruktur zum zentralen Baustein der TypeScript-Lösung
wurde.


\subsection{Einordnung für das Projekt und Zwischenfazit}
\label{subsec:rs-api-einordnung-fazit}

Mit \mintinline{text}{RecursiveSet} wurde in diesem Kapitel nicht einfach eine
weitere Collection-Klasse eingeführt. Die Datenstruktur erfüllt vielmehr eine
zentrale Funktion innerhalb der gesamten Portierung: Sie stellt genau das
Mengenverhalten bereit, das in den ursprünglichen Python-Notebooks fachlich
vorausgesetzt wurde, im eingebauten JavaScript-\mintinline{text}{Set} für
zusammengesetzte Werte jedoch nicht in derselben Form verfügbar ist.

Gerade darin liegt ihre Bedeutung für das Projekt. Die Portierung sollte nicht
nur einzelne Python-Programme syntaktisch nach TypeScript übertragen, sondern
die zugrunde liegende Denkweise erhalten. In vielen Notebook-Beispielen wird
mit Mengen von Zuständen, Paaren, Teilmengen oder anderen zusammengesetzten
Werten gearbeitet. Damit solche Konstruktionen in TypeScript ebenso natürlich
formulierbar bleiben, musste eine Datenstruktur bereitgestellt werden, die
nicht bloß Referenzen verwaltet, sondern Werte nach ihrer inhaltlichen
Struktur behandelt.

Dieses Kapitel hat deshalb zunächst die Benutzung von
\mintinline{text}{RecursiveSet} aus Anwendersicht aufgebaut. Ausgehend vom
Mengenproblem beim Übergang von Python nach TypeScript wurde gezeigt, wie die
Bibliothek grundlegende Mengenoperationen, strukturbezogene Mitgliedschaft,
kartesische Produkte sowie funktionale Transformationen auf Mengen
bereitstellt. Dadurch wurde sichtbar, dass \mintinline{text}{RecursiveSet}
nicht nur einzelne Hilfsmethoden bündelt, sondern eine konsistente
Arbeitsoberfläche für die Mengenmodellierung schafft.

Zugleich wurde deutlich, dass diese Arbeitsoberfläche nicht voraussetzungslos
ist. Die Bibliothek erlaubt nicht beliebige Werte, sondern verlangt bei
komplexeren Strukturen einen klaren Vertrag über Gleichheit, Hashing und
Darstellung. Auch der Lebenszyklus zwischen veränderlicher und eingefrorener
Menge gehört zu dieser bewussten API-Gestaltung. Gerade diese Regeln machen
die Datenstruktur für das Projekt fachlich belastbar, statt sie nur bequem
wirken zu lassen.

Für die Studienarbeit ergibt sich daraus eine klare Einordnung:
\mintinline{text}{RecursiveSet} ist die zentrale Kompensation des
Mengenproblems in der TypeScript-Portierung. Erst durch diese Datenstruktur
konnte die aus den Python-Notebooks bekannte Arbeit mit Mengen nicht nur
oberflächlich, sondern auch semantisch überzeugend übernommen werden. Die
Bibliothek bildet damit eine wesentliche Brücke zwischen der mathematischen
Denkweise der Lehrmaterialien und ihrer technischen Realisierung in
TypeScript.

Damit ist zugleich die Grenze dieses Kapitels erreicht. Bisher stand im
Vordergrund, \emph{welches} Verhalten \mintinline{text}{RecursiveSet}
bereitstellt und \emph{warum} dieses Verhalten für das Projekt notwendig ist.
Im folgenden Kapitel wechselt die Perspektive nun von der Benutzung zur
internen Konstruktion. Dort wird Schritt für Schritt untersucht, wie Hashing,
strukturelle Gleichheit, Speicherlayout und die Algorithmen für Einfügen,
Suchen und Löschen dieses Verhalten technisch umsetzen.

